%% ma: L12-B08-0071
%% lop: 12
%% chuong: 2
%% bai: 8
%% dang: 12.8.32
%% loai: TLN
%% muc: TH
%% dap_an: "2300"
%% nguon: "LaTeX giáo viên gửi 21/10/2026 (M05, Vectơ và hệ trục tọa độ trong không gian, 107 câu), Phần III Câu 11"
%% kien_thuc: "Bài 8 (tọa độ vectơ, tích vô hướng), ứng dụng tính công"
%% trang_thai: chinh-thuc
%% ghi_chu: "Giáo viên đã duyệt 21/10/2026. Phần III Câu 20 của tài liệu (cần cẩu kéo khối bê tông) chỉ đổi số so với câu này nên đã bỏ"
\begin{ex}
Một thiết bị kéo hàng trong nhà máy tác dụng một lực $\overrightarrow{F}=(200;300;-100)$ (đơn vị: N) để làm dịch chuyển một toa hàng từ vị trí $M(1;2;0)$ đến vị trí $N(6;7;2)$ (đơn vị tọa độ: mét). Tính công $A$ (đơn vị: J) sinh ra bởi lực $\overrightarrow{F}$ trong quá trình dịch chuyển đó.
\shortans{2300}
\loigiai{$\overrightarrow{MN}=(5;5;2)$, $A=\overrightarrow{F}\cdot\overrightarrow{MN}=1000+1500-200=2300$ (J).}
\end{ex}
